\section{Introduction}

Distributed Point Functions (DPFs) have emerged as a core primitive for function secret sharing\cite{fss,dpf2}, enabling two parties to hold compact keys that expand into additive shares of a unit vector of length~$N$ \cite{dpf,dpf2}. This remarkable ability to compress a sparse vector into two short keys has found diverse applications: pseudorandom correlation generators (PCGs) for secure computation \cite{pcg2019,boyle2020efficient,pcgZ2k,anyField}, efficient two-server private information retrieval (PIR) \cite{demmler2018pir,inc-offline-online-pir}, private set intersection (PSI) \cite{blazingPSI,volePSI,mpPSI}, oblivious RAM \cite{dpfKeygen}, and federated analytics \cite{prio,vdaf}.

Many such applications require not just a single hidden position, but several positions at once. For example, Ring-LPN based PCGs \cite{boyle2020efficient} or the recent FOLEAGE protocol \cite{foleage} rely on compact encodings of sparse structures with Hamming weight $t$, while batched PIR queries and PSI protocols involve sets of multiple elements. A natural generalization is the \emph{Distributed Multi-Point Function} (DMPF), which shares a vector with up to $t$ non-zero entries. Naïvely, one can run $t$ independent DPFs, but this inflates cost by a factor of $t$, yielding $O(Nt)$ work instead of the $O(N)$ ideal.

Prior DMPF approaches highlighted two main obstacles. First, efficient \emph{distributed} key generation is difficult: early proposals often assumed that one party learns all indices, or that a trusted dealer provides keys \cite{dmpf}. Second, optimized cuckoo-hash based schemes \cite{cuckooDmpf} reduced asymptotic overhead but introduced iterative insertion and eviction subroutines. When the indices are secret-shared, such procedures become expensive MPC protocols that dominate concrete performance.

To avoid confusion about what is conceptually new versus what is fixed, we emphasize that while prior work established DMPFs and cuckoo-based constructions in settings where indices are known or a dealer is available, no prior work provided an efficient \emph{fully distributed} DMPF suitable for MPC. Our contribution is unlocking practical performance: we eliminate iterative insertion, avoid heavy MPC subroutines, and give the first practical distributed DMPF key generation with near-linear cost.

\paragraph{Our contribution.}
We present new, \emph{fully distributed} DMPF constructions that close this gap and enable scalable, MPC-friendly key generation. Our central tool is the \emph{Reverse Cuckoo} paradigm: instead of fixing hash functions first and searching for placements, we randomly assign items to buckets and then solve for hash functions consistent with this assignment. This inversion eliminates iterative eviction loops and yields a one-shot linear-algebraic subroutine that integrates smoothly into MPC. We also develop complementary bucketing and big-state approaches, and show how all these constructions integrate into PCGs from Ring-LPN and Stationary LPN.

\emph{Performance.} Our implementation achieves multi-million $\mathbb F_q$-OLE/sec online throughput with modest, $n$-independent per-expansion communication. For the 64 bit ``Goldilocks'' prime\cite{hamburg2015ed448,goldilocks}, $q=2^{64}-2^{32}-1$, generates rates climb from $0.89{\times}10^6$ to $1.58{\times}10^6$ to $1.81{\times}10^6$ OLEs/s as $n$ grows from $2^{16}$ to $2^{20}$, with per-expansion communication $\approx13$\,MB and setup times $2.63$-$12.49$\,s and $149$-$170$\,MB setup communication. For a 31-bit prime \emph{Fp31} is $1.4$-$1.5\times$ faster online (e.g., $2.71{\times}10^6$ ops/s at $n{=}2^{20}$) and uses slightly less setup bandwidth. Comparing our DMPF with the status-quo “sum of $t$ DPFs,” we see over an order-of-magnitude higher throughput (e.g., $885{,}621$ vs.\ $37{,}686$ OLEs/s at $n{=}2^{16}$, Goldilocks). 

Compared to \cite{mp-spdz}, the state of art for Beaver triple generation based on FHE, which achieves $275{,}548$-$309{,}588$ OLEs/s at $n{=}2^{20}$, our online throughput is $6$-$10\times$ higher with substantially lower per-expansion communication. Moreover, the recent PCG based implementation of \cite{silentium} using sum of DPFs achieves an lower throughput of just 22,800 OLEs per second, although on different hardware and larger prime.  In conclusion, generating Beaver triples is now an order of magnitude more efficient. See \sectionref{sec:eval} for a detailed comparison. 
